\section{LDA 与收缩：判别方向的几何与高维失效}

上一节把 LDA 写成共享协方差的 GDA——Bayes 规则下二次项抵消，边界是线性的。这回答了``LDA 从哪里来''，却没有回答两个更实际的问题：这个线性边界在几何上是什么？为什么高维小样本时它常常失效，以及收缩怎样修复。这两个问题分别通向 Fisher 的原始视角与 Part 11 的随机矩阵理论。

\subsection{Fisher 视角：让类间分离与类内弥散之比最大}

Fisher 的出发点不是概率模型，而是一个可分性准则：把数据投影到方向 $w$ 上，投影后两类均值之差应尽可能大，类内方差应尽可能小。写成比值，

\[
J(w)=\frac{(w^\trans(\mu_1-\mu_0))^2}{w^\trans\Sigma w}
=\frac{w^\trans S_b w}{w^\trans S_w w},
\]

其中 $S_b=(\mu_1-\mu_0)(\mu_1-\mu_0)^\trans$ 是类间散度（二分类时秩为 1），$S_w=\Sigma$ 是共享的类内散度。这是 Rayleigh 商；对 $w$ 求极值，一阶条件给出

\[
S_b w=\lambda S_w w,
\]

而 $S_b w$ 总是 $(\mu_1-\mu_0)$ 的倍数，所以解就是

\[
w^\ast\propto S_w^{-1}(\mu_1-\mu_0)=\Sigma^{-1}(\mu_1-\mu_0).
\]

与 GDA 的判别函数 $\delta_k(x)$ 对比：那里的线性系数正是 $\Sigma^{-1}\mu_k$。两条路线——Bayes 规则下的生成式分类、Rayleigh 商下的判别方向——给出同一个方向。这不是巧合：共享协方差的 Gaussian 是让 Fisher 准则等于 Bayes 最优的那个分布族。离开 Gaussian，Fisher 方向仍是一个合理的线性投影，但不再享有``最优''的称号；这是它作为判别式方法与 GDA 作为生成式方法的分工。

几何上，$\Sigma^{-1}$ 做的事是``白化''：先用 $\Sigma^{-1/2}$ 把椭球形类内云变成球形，再在球面几何里找均值连线方向。类内弥散大的方向被压缩，类间差异因此按``信噪比''而非欧氏距离度量——这正是 Mahalanobis 距离的含义（Part 03 度量与 Part 11 马氏距离在此处汇合）。

\subsection{高维失效：$\widehat\Sigma^{-1}$ 是失真最厉害的环节}

实际中 $\Sigma$ 未知，用样本协方差 $\widehat\Sigma$ 代入。Part 11 的随机矩阵理论已经给出警告：当 $p/n\to c>0$，$\widehat\Sigma$ 的特征值弥散到 MP 律的支撑 $[(1-\sqrt c)^2,(1+\sqrt c)^2]$ 上，即使真实 $\Sigma=I$。求逆会放大这种失真：最小的样本特征值靠近 $(1-\sqrt c)^2$，其倒数把对应方向上的噪声放大到 $1/(1-\sqrt c)^2$ 倍。$c\to1$ 时这个因子爆炸；$p\ge n$ 时 $\widehat\Sigma$ 干脆奇异，$w^\ast$ 根本不存在。

判别方向 $w^\ast$ 对 $\Sigma^{-1}$ 的依赖是线性的，所以谱失真直接传递为方向失真：$\widehat w$ 会在``样本协方差的小特征值方向''上获得虚假的大分量——那些方向恰恰是噪声主导的方向。表现上，训练集里分离完美的方向在新数据上崩溃，这是高维 LDA 的经典失败模式，也是``模型容量''叙事之外一个更精确的解释：问题不在线性边界这个假设，而在 $\widehat\Sigma^{-1}$ 这个估计量。

\subsection{收缩 LDA：向结构化目标折回}

修复沿用 Part 11 的两条思路，而且在这里有明确的统计含义：

\textbf{Ledoit--Wolf 收缩。}把 $\widehat\Sigma$ 向一个结构化目标（如单位阵的倍数 $\bar\sigma^2 I$）收缩：

\[
\widetilde\Sigma=(1-\rho)\widehat\Sigma+\rho\bar\sigma^2 I,
\qquad \rho\in[0,1].
\]

$\rho$ 由最小化 $\widetilde\Sigma$ 与真实 $\Sigma$ 的期望误差（MSE）选出，有闭式估计。几何上，收缩把 MP 律弥散开的特征值向中间挤压：大特征值压小、小特征值抬大，求逆时的噪声放大被限制在 $1/((1-\rho)(1-\sqrt c)^2+\rho\bar\sigma^2)$ 量级。统计上，这是偏差--方差权衡的又一个实例：收缩引入对 $\Sigma$ 的偏差假设（各方向方差相近），换取 $\widetilde\Sigma^{-1}$ 的方差大幅下降。Ridge 回归中$\lambda$ 的角色、James--Stein 估计中收缩因子的角色，与这里的 $\rho$ 是同构的。

\textbf{对角收缩与正则化 LDA。}更强的假设是直接取对角阵 $\widetilde\Sigma=\operatorname{diag}(\widehat\Sigma)$，即假设类内各特征不相关——这退化为 Naive Bayes 的 Gaussian 版本。看似倒退，在高维却常常优于全协方差 LDA：估计参数从 $O(p^2)$ 降到 $O(p)$，相关性误差小于求逆误差。两端的连续谱——全协方差、收缩、对角——就是``模型复杂度随样本量自适应''的一个干净例子。

多分类推广（$K>2$）：$S_b$ 的秩至多是 $K-1$，Fisher 判别给出一组方向（$S_w^{-1}S_b$ 的特征向量），构成降到 $K-1$ 维的判别子空间。高维下的处理与二分类相同：先收缩 $S_w$，再做广义特征分解。

\subsection{边界清单}

\begin{itemize}
\tightlist
\item
  LDA 的两个前提——类条件 Gaussian、共享协方差——缺一不可验。各类 $\Sigma_k$ 差异大时应改用 QDA，且 QDA 对协方差估计误差更敏感，高维下更需收缩。
\item
  Fisher 方向在 Gaussian 下最优，在非 Gaussian 下只是一个投影方向；若类条件分布重尾或不对称，稳健散度（中位数、 trimming）比样本均值协方差更可靠。
\item
  收缩强度 $\rho$ 本身也由数据估计，自带方差；小样本时交叉验证选 $\rho$ 常优于闭式 MSE 公式。
\item
  收缩修复的是$\widehat\Sigma$的谱，不修复模型错设。类边界本质上非线性时，收缩 LDA 与朴素 LDA 一样无能为力——那是核方法或深模型的领域（Part 13 核方法单元）。
\end{itemize}

\subsection{落点}

LDA 把三件事接成一条链：生成式建模（GDA 共享协方差）给出判别函数，几何准则（Fisher 的 Rayleigh 商）给出同一方向的判别式推导，高维统计（MP 律与收缩）解释并修复它的失效。这条链的两个端点——Part 13 的分类与 Part 11 的随机矩阵——原本分属``机器学习''与``数理统计''，LDA 是它们最自然的交汇点之一。
